\chapter{The Dynamic Window Approach}
\label{ch:ch14}
% Function names for pseudocode are prefixed with "Dwa" to stay unique
% across the book.
\SetKwFunction{DwaCommand}{DwaCommand}
\SetKwFunction{DwaCandidates}{WindowCandidates}
\SetKwFunction{DwaFreeDistance}{FreeDistance}
\SetKwFunction{DwaLoop}{DwaControlLoop}
\SetKwFunction{DwaLookahead}{LookaheadPoint}

\begin{objectives}
\begin{itemize}
  \item Explain what the dynamic window approach (DWA) computes: the best velocity command among those the robot can actually reach in the next control interval, judged by short simulated rollouts.
  \item Construct the three velocity sets of DWA (possible, admissible, dynamic window) for a differential-drive robot and for a velocity-controlled drone, and derive the braking condition $v \le \sqrt{2\,\dist\,a_b}$ together with its discrete-time refinement.
  \item Evaluate the heading, clearance and velocity terms of the objective by hand, explain why they must be normalised, and predict what happens when each weight dominates.
  \item Prove that an admissible command can always be followed by a safe one, and show by example that DWA is a local method that can get stuck in a U-shaped trap and can oscillate.
  \item Implement DWA in \python for a drone in a continuous 2D or 3D world, connect its heading term to a global path, and create and repair its failure cases.
\end{itemize}
\end{objectives}

% ---------------------------------------------------------------------
\section{Why this matters for a drone swarm}
\label{sec:ch14-motivation}

The local layer of the hybrid architecture in \cref{ch:ch24} has one job: when
something unexpected enters the safety horizon of a drone, produce a velocity
for the next fraction of a second that is safe \emph{and} keeps the drone
roughly on its nominal route. \Cref{ch:ch12,ch:ch13} solved this with velocity
obstacles and \orca. Both treat the drone as a point that can change its
velocity instantly: \orca picks any velocity outside a set of forbidden
half-planes, and the drone is assumed to adopt it at once. For a light
quadrotor at moderate speed this is a fair approximation. For a heavy
delivery drone, a fixed-wing aircraft or any vehicle whose acceleration is
limited to a fraction of what \orca asks for, it is not: the drone is told to
fly at a velocity it cannot reach within the control interval, and the
guarantee of \cref{ch:ch13} evaporates.

The dynamic window approach\index{dynamic window approach|see{DWA}}
(\textbf{DWA})\index{DWA} of Fox, Burgard and Thrun~\cite{fox1997dwa} starts
from the opposite end. It asks first which velocities the robot can reach in
the next $\dt$ seconds given its acceleration limits, then throws away those
from which it could no longer stop before the nearest obstacle, simulates the
short trajectory that each remaining candidate would produce, scores the
candidates by how well they head towards the goal, how much clearance they
keep and how fast they are, and executes the winner for one interval. Then it
starts over. DWA was developed for the wheeled robot RHINO, which had to drive
at up to about a metre per second through crowded rooms with a laser scanner
and sonar as its only knowledge of the world, and it has been the default local
planner of many mobile-robot navigation stacks ever since.

For the drone system this gives a second local layer, an alternative to \orca
when the drone's dynamics matter more than reciprocity with the other agent:
DWA never asks for a velocity the drone cannot produce, it works against any
obstacle that can be sensed or predicted, and it does not assume that the
other party shares the effort. The price is that it is a greedy, memoryless
method. It can drive into a dead end and stay there, and it can oscillate
between two equally attractive commands. The Week~7 milestone of the study
plan (\cref{ch:appA}) is precisely to know when such local methods work and
when they fail, and this chapter and the next (\cref{ch:ch15}) are built
around that question.

The chapter proceeds as follows. \Cref{sec:ch14-intuition} gives the idea.
\Cref{sec:ch14-problem} defines the velocity space, the rollouts, the three
restrictions and the objective, in the differential-drive form of the original
paper and in the drone form used in the rest of the chapter.
\Cref{sec:ch14-algorithm} gives the pseudocode and \cref{sec:ch14-example}
runs it by hand on a drone approaching an obstacle, with every number taken
from \code{ch14\_dwa.py}. \Cref{sec:ch14-properties} proves what DWA
guarantees (safety with respect to static obstacles, dynamic feasibility) and
what it does not (completeness), \cref{sec:ch14-tuning} studies the weights
in a generated experiment, \cref{sec:ch14-variants} collects the remedies and
compares DWA with \orca and potential fields, and
\cref{sec:ch14-implementation,sec:ch14-drone} cover the implementation and
the place of DWA in the drone system.

% ---------------------------------------------------------------------
\section{The idea in plain words}
\label{sec:ch14-intuition}

Think of the problem from the point of view of the flight controller rather
than the map. Every $\dt = 0.25$ seconds it must hand a velocity command to the
motors. Whatever the global planner says, the drone is right now flying at
some velocity $\vel_a$, and in a quarter of a second it can change that
velocity only by $a_{\max}\dt$ along each axis. The set of commands it can
actually execute is therefore a small box around $\vel_a$ in velocity space,
the \textbf{dynamic window}.\index{dynamic window} Everything outside the
window is wishful thinking.

Inside the window DWA reasons like a careful driver. For each candidate
velocity it imagines flying with that velocity for a short while, a
\textbf{rollout},\index{rollout} and measures how far it could go before
hitting something. If that distance is shorter than the distance needed to
brake to a standstill, the candidate is thrown away: choosing it would commit
the drone to a collision even if the very next command were full braking.
This is the \textbf{admissibility} test.\index{admissible velocity} The
remaining candidates are all safe for at least one more step, and among them
DWA simply picks the one with the best score, a weighted sum of three
ingredients: does the rollout head towards the goal, how much free space does
it have ahead, and how fast is it. The winner is sent to the motors, the
drone moves, the sensors deliver a new picture of the obstacles, and the whole
computation is repeated from the new velocity. \Cref{fig:ch14-pipeline}
shows the loop.

\begin{figure}[tb]
  \centering
  \inputfigure{ch14/pipeline}
  \caption{One control step of DWA. The current velocity $\vel_a$ spans the
  dynamic window; the candidates in it are rolled out against the obstacles;
  those that cannot brake in time are discarded; the rest are scored by
  heading, clearance and velocity; the best one is executed for $\dt$ and the
  loop repeats. Nothing is remembered between steps except the velocity
  itself.}
  \label{fig:ch14-pipeline}
\end{figure}

Two things distinguish this from the velocity-obstacle methods of
\cref{ch:ch12,ch:ch13}. First, DWA searches the space of \emph{achievable}
commands over one control interval, so every command it returns respects the
robot's dynamics by construction; \orca searches all velocities and hopes the
robot can follow. Second, DWA judges a command by simulating it: the rollout
is an explicit prediction of where the robot will be, along a circular arc for
a wheeled robot and along a straight segment for a velocity-controlled drone,
and any obstacle representation that supports a ``how far is it free in this
direction'' query can be plugged in. The velocity obstacle, in contrast, is an
analytic construction for discs moving at constant velocity.

\begin{keyidea}
Search only the velocities you can reach in the next control interval, keep
only those from which you can still stop before the nearest obstacle, simulate
each for a short horizon, and execute the one whose rollout best combines
heading towards the goal, clearance and speed. Safety comes from the braking
test; progress comes from the score; memory comes from nowhere, which is why
DWA is fast, respects dynamics, and can still get stuck.
\end{keyidea}

% ---------------------------------------------------------------------
\section{Problem statement and notation}
\label{sec:ch14-problem}

\subsection{The velocity space}
\label{sec:ch14-velocity-space}

DWA plans in \textbf{velocity space}:\index{velocity space} the coordinates of
the search are not positions but the components of the velocity command that
is held constant during the next control interval $\dt$. Two robot models are
used in this chapter.

\begin{definition}[Velocity space of a differential-drive robot]
\label{def:ch14-diffdrive}
A \textbf{differential-drive robot}\index{differential drive} has the pose
$(x, y, \theta)$ and is commanded by a translational velocity $v$ and a
rotational velocity $\omega$. Its velocity space is the set of pairs
$(v, \omega)$. During an interval in which $(v, \omega)$ is constant the robot
moves on a circular arc of radius $v/\omega$ (a straight line if
$\omega = 0$), so that after time $t$
\begin{equation}
  x(t) = x + \frac{v}{\omega}\bigl(\sin(\theta + \omega t) - \sin\theta\bigr),\quad
  y(t) = y - \frac{v}{\omega}\bigl(\cos(\theta + \omega t) - \cos\theta\bigr),\quad
  \theta(t) = \theta + \omega t.
  \label{eq:ch14-arc}
\end{equation}
\end{definition}

\begin{definition}[Velocity space of a velocity-controlled drone]
\label{def:ch14-drone}
A \textbf{velocity-controlled drone}\index{drone!velocity-controlled} is the
point-mass model of \cref{ch:ch02}: its state is the position $\pos \in \R^3$
and its command is the velocity $\vel = (v_x, v_y, v_z)$, held constant for
$\dt$, so that $\pos(t) = \pos + \vel\,t$ for $0 \le t \le \dt$. The velocity
space is $\R^3$ (or $\R^2$ for flight at constant altitude). The drone has a
speed limit $v_{\max}$ and a per-axis acceleration limit $a_{\max}$: two
consecutive commands $\vel_a$ and $\vel$ must satisfy
$\abs{v_k - v_{a,k}} \le a_{\max}\dt$ for every axis $k$.
\end{definition}

The drone is \textbf{holonomic}:\index{holonomic} it can move in any direction
regardless of its orientation, so a rollout is a straight segment
(\cref{fig:ch14-rollouts}, right). The wheeled robot is not: it can only move
along its heading, and a rollout is an arc (\cref{fig:ch14-rollouts}, left).
Everything else in DWA is the same for both models. We use the drone form from
\cref{sec:ch14-algorithm} on and write $\vel_a$ for the velocity the drone is
currently flying with, exactly as Fox et al.\ write $(v_a, \omega_a)$ for the
actual velocity of their robot.

\begin{figure}[tb]
  \centering
  \inputfigure{ch14/rollouts}
  \caption{Rollouts are the trajectories the robot would follow if it kept a
  candidate command for a while. Left: a differential-drive robot heading
  along $x$ at $v = 1$~m/s; the arcs are $\omega = 0, \pm 0.5, \pm 1$~rad/s,
  with radius $v/\omega$. The free distance $\dist(v, \omega)$ is measured
  along the arc up to the first contact with the inflated obstacle. Right: a
  velocity-controlled drone; every rollout is a straight segment
  $\pos + \vel\,t$, and $\dist(\vel)$ is a ray cast in the direction of
  $\vel$. Dotted rollouts hit the obstacle within the sensing range
  $d_{\max}$.}
  \label{fig:ch14-rollouts}
\end{figure}

Obstacles are handled in the inflated configuration space of \cref{ch:ch02}:
every obstacle is grown by the drone's radius $r$ and the drone becomes a
point. For a candidate $\vel \ne \vect{0}$ the \textbf{free distance}
\begin{equation}
  \dist(\vel) = \min\Bigl(d_{\max},\ \text{distance from $\pos$ along the ray in direction $\vel/\norm{\vel}$ to the first inflated obstacle}\Bigr)
  \label{eq:ch14-dist}
\end{equation}
is what the rollout measures.\index{free distance} The cap $d_{\max}$ is the
sensing range, or simply the distance beyond which we stop caring. For
$\vel = \vect{0}$ we set $\dist(\vect{0}) = d_{\max}$: standing still never
runs into a static obstacle. For discs the ray cast is the ray--circle
intersection of \cref{ch:ch02}; for axis-aligned boxes it is a slab test; for
a point cloud or an occupancy grid it is a march along the ray.

\subsection{The three restrictions}
\label{sec:ch14-restrictions}

The search set of DWA is the intersection of three sets in velocity space
(\cref{fig:ch14-spaces}). We state each for the differential drive, as in
\cite{fox1997dwa}, and for the drone.

\begin{figure}[tb]
  \centering
  \begin{subfigure}[b]{0.32\textwidth}\centering
    \inputfigure{ch14/space-possible}
    \caption{possible velocities $V_s$}
  \end{subfigure}\hfill
  \begin{subfigure}[b]{0.32\textwidth}\centering
    \inputfigure{ch14/space-admissible}
    \caption{admissible velocities $V_a$}
  \end{subfigure}\hfill
  \begin{subfigure}[b]{0.32\textwidth}\centering
    \inputfigure{ch14/space-window}
    \caption{dynamic window $V_d$}
  \end{subfigure}
  \caption{The three restrictions in the $(\omega, v)$ space of a
  differential-drive robot. (a) The velocity limits give a rectangle.
  (b) Velocities from which the robot could not stop before the obstacle on
  its arc are cut away (hatched); the bite is deepest for the curvatures that
  lead straight into the obstacle. (c) The dynamic window is the rectangle of
  velocities reachable from $(v_a, \omega_a)$ within one interval; the search
  set $V_r$ (dark) is its intersection with the first two sets. Only the grid
  points inside $V_r$ are evaluated.}
  \label{fig:ch14-spaces}
\end{figure}

\begin{definition}[Possible velocities]
\label{def:ch14-possible}
The set of \textbf{possible velocities}\index{DWA!possible velocities} $V_s$
contains the commands the actuators can produce at all. For the differential
drive, $V_s = [0, v_{\max}] \times [-\omega_{\max}, \omega_{\max}]$ (or
$[v_{\min}, v_{\max}]$ if driving backwards is allowed). For the drone,
$V_s = \set{\vel : \norm{\vel} \le v_{\max}}$, a ball, or a box if the axes
have separate limits.
\end{definition}

\begin{definition}[Admissible velocities]
\label{def:ch14-admissible}
A velocity is \textbf{admissible}\index{admissible velocity}\index{DWA!admissible velocities}
if the robot can still come to a standstill before it reaches the nearest
obstacle on its rollout. With the braking decelerations $\dot v_b$ and
$\dot\omega_b$ of the differential drive,
\begin{equation}
  V_a = \set{(v, \omega) \given v \le \sqrt{2\,\dist(v,\omega)\,\dot v_b}
  \ \text{ and }\ \omega \le \sqrt{2\,\dist(v,\omega)\,\dot\omega_b}}.
  \label{eq:ch14-Va}
\end{equation}
For the drone with braking deceleration $a_b$,
$V_a = \set{\vel \given \norm{\vel} \le \sqrt{2\,\dist(\vel)\,a_b}}$.
\end{definition}

\paragraph{Where the square root comes from.}
A robot that moves at speed $v$ and decelerates at the constant rate $a_b$ has
speed $v - a_b t$ after $t$ seconds and stops at $t_s = v/a_b$. The distance it
covers meanwhile is
\begin{equation}
  s_b(v) = \int_0^{t_s} (v - a_b t)\,dt = v\,t_s - \tfrac{1}{2} a_b t_s^2
         = \frac{v^2}{2 a_b},
  \label{eq:ch14-braking}
\end{equation}
the \textbf{braking distance}.\index{braking distance} The robot is safe if the
braking distance fits into the free distance, $v^2/(2a_b) \le \dist$, and
solving for $v$ gives $v \le \sqrt{2\,\dist\,a_b}$. The condition on
$\omega$ in \cref{eq:ch14-Va} is the same argument applied to the rotation, so
that the robot can also stop turning; Fox et al.\ measure $\dist$ along the
arc of the curvature $\omega/v$, which is why it is written $\dist(v,\omega)$.

\paragraph{The discrete-time refinement.}
\Cref{eq:ch14-braking} assumes that braking starts immediately. In a
digital control loop the chosen command is held for the whole interval $\dt$
before the next command, at best full braking, can begin. The distance to
standstill is then at most
\begin{equation}
  s_b^{\dt}(v) = v\,\dt + \frac{v^2}{2 a_b},
  \label{eq:ch14-braking-discrete}
\end{equation}
and the admissibility test becomes $s_b^{\dt}(\norm{\vel}) \le \dist(\vel)$,
or, solving the quadratic for the speed,
\begin{equation}
  \norm{\vel} \le v_{\mathrm{adm}}(\dist) = -a_b\dt + \sqrt{a_b^2\dt^2 + 2\,a_b\,\dist}.
  \label{eq:ch14-vadm}
\end{equation}
For $\dt \to 0$ this is the condition of \cref{eq:ch14-Va}. For the drone of
\cref{sec:ch14-example}, with $a_b = 2$~m/s$^2$ and $\dt = 0.25$~s at
$0.608$~m from an obstacle, \cref{eq:ch14-Va} allows $1.56$~m/s while
\cref{eq:ch14-vadm} allows $1.14$~m/s. The difference is not academic: the
first version of \code{ch14\_dwa.py} used \cref{eq:ch14-Va} and drove the
drone into the obstacle within a few steps. \Cref{sec:ch14-properties}
proves that \cref{eq:ch14-vadm} is enough, and all code and numbers in this
chapter use it. (The extra term is exactly the pitfall ``a dynamic window so
small that the robot cannot react'' seen from the braking side: a command is
only as safe as the reaction it still allows.)

\begin{definition}[Dynamic window]
\label{def:ch14-window}
The \textbf{dynamic window}\index{dynamic window}\index{DWA!dynamic window}
$V_d$ contains the velocities reachable from the current velocity within one
control interval under the acceleration limits. For the differential drive
with the maximal accelerations $\dot v_{\max}$ and $\dot\omega_{\max}$,
\begin{equation}
  V_d = \set{(v, \omega) \given
    v \in [v_a - \dot v_{\max}\dt,\ v_a + \dot v_{\max}\dt]
    \ \text{ and }\
    \omega \in [\omega_a - \dot\omega_{\max}\dt,\ \omega_a + \dot\omega_{\max}\dt]}.
  \label{eq:ch14-Vd}
\end{equation}
For the drone with the per-axis limit $a_{\max}$,
$V_d = \set{\vel \given \abs{v_k - v_{a,k}} \le a_{\max}\dt \text{ for all } k}$,
a cube of side $2 a_{\max}\dt$ centred on $\vel_a$.
\end{definition}

\begin{definition}[Search set]
\label{def:ch14-search-set}
The \textbf{search set}\index{DWA!search set} of one DWA step is
$V_r = V_s \cap V_a \cap V_d$. In practice it is discretised: the window is
covered by a grid of resolution $\Delta v$ centred on $\vel_a$, and only the
grid points in $V_r$ are evaluated. With $n = a_{\max}\dt/\Delta v$ the grid
has $(2n+1)^{\!\dim}$ points before the admissibility test.
\end{definition}

The three sets play different roles. $V_s$ is fixed. $V_d$ moves with the
robot's own velocity and is small: with $a_{\max} = 2$~m/s$^2$ and
$\dt = 0.25$~s it spans only $\pm 0.5$~m/s around $\vel_a$, so the robot can
change course only gradually and the dynamics are respected without any
further modelling. $V_a$ moves with the obstacles and is where all the
geometry enters. Note that the window is a constraint on the \emph{change} of
velocity, not on the velocity, which is why a fast robot approaching a wall
needs several steps to slow down and why the braking test must look further
ahead than one step.

\subsection{The objective}
\label{sec:ch14-objective}

Among the velocities in $V_r$, DWA maximises
\begin{equation}
  G(\vel) = \sigma\bigl(\alpha \cdot \mathrm{heading}(\vel)
             + \beta \cdot \mathrm{clearance}(\vel)
             + \gamma \cdot \mathrm{velocity}(\vel)\bigr),
  \label{eq:ch14-G}
\end{equation}
where $\alpha, \beta, \gamma \ge 0$ are weights and $\sigma$ is a smoothing
discussed below.\index{DWA!objective function} Fox et al.\ call the three
terms \emph{target heading}, \emph{clearance} (written $\mathrm{dist}$) and
\emph{velocity}. In the normalised form used in this chapter all three lie
in $[0, 1]$.

\begin{definition}[Heading term]
\label{def:ch14-heading}
Let $\pos' = \pos + \vel\,\dt$ be the position predicted after the next
interval and $\theta(\vel) \in [0, \pi]$ the angle between $\vel$ and the
direction $\vect{t} - \pos'$ from $\pos'$ to the target point $\vect{t}$
(the goal, or a point on the global path, \cref{sec:ch14-variants}). The
\textbf{heading term}\index{heading term}\index{DWA!heading term} is
\begin{equation}
  \mathrm{heading}(\vel) = 1 - \frac{\theta(\vel)}{\pi},
  \label{eq:ch14-heading}
\end{equation}
with $\mathrm{heading}(\vect{0}) = 0$. It is $1$ when the rollout points
straight at the target and $0$ when it points away. Fox et al.\ use
$180^\circ - \theta$ with $\theta$ measured at the position and orientation
the robot is predicted to reach after the next interval; for the
differential drive the orientation changes along the arc, so this predicted
pose matters more than for the drone.
\end{definition}

\begin{definition}[Clearance term]
\label{def:ch14-clearance}
The \textbf{clearance term}\index{clearance term}\index{DWA!clearance term} is
the normalised free distance of the rollout,
\begin{equation}
  \mathrm{clearance}(\vel) = \frac{\min(\dist(\vel), d_{\max})}{d_{\max}} \in [0, 1].
  \label{eq:ch14-clearance}
\end{equation}
Fox et al.\ use the unnormalised $\dist(v, \omega)$ and set it to a large
constant when no obstacle lies on the curvature.
\end{definition}

\begin{definition}[Velocity term]
\label{def:ch14-velocity}
The \textbf{velocity term}\index{velocity term}\index{DWA!velocity term}
rewards speed,
\begin{equation}
  \mathrm{velocity}(\vel) = \frac{\norm{\vel}}{v_{\max}} \in [0, 1].
  \label{eq:ch14-velocity}
\end{equation}
For the differential drive it is the translational speed $v$; the rotation
$\omega$ is not rewarded.
\end{definition}

\paragraph{Why the terms are normalised.}
The three terms measure different things: an angle, a distance and a speed. In
their raw units a weighted sum has no meaning: the same weights that balance
heading against clearance in a room measured in metres would let clearance
dominate everything if the map were in centimetres, and would have to be
retuned for a drone that flies ten times faster. Dividing each term by its
maximal value ($\pi$, $d_{\max}$, $v_{\max}$) makes the weights dimensionless
fractions that say how much each concern counts, and lets the same
$(\alpha, \beta, \gamma)$ be reused across scenes and vehicles.\index{normalisation}
Many implementations go one step further and normalise each term over the
current candidate set, so that the best candidate of each term scores $1$;
this makes the weights even more portable but also makes $G$ depend on which
candidates happen to be admissible.

\paragraph{The smoothing $\sigma$.}
The original objective is smoothed: $\sigma$ averages the weighted sum over
the neighbouring cells of the velocity grid. Its purpose is to prefer
commands whose neighbours are also good, which pushes the robot away from
narrow ``good'' spots at the edge of the admissible region and gives more
side clearance to obstacles.\index{smoothing} A candidate whose neighbours are
inadmissible (score $0$) is pulled down. \Cref{sec:ch14-variants} relates
$\sigma$ to hysteresis; the code offers it as an option, and the worked
example and the experiments use $\sigma = \mathrm{id}$, that is, no
smoothing, to keep the numbers transparent.

\begin{table}[tb]
  \centering\small
  \caption{Symbols of this chapter and the values used in \code{ch14\_dwa.py}
  (class \code{DwaParams}). The drone flies at constant altitude in the
  examples, so velocities have two components.}
  \label{tab:ch14-parameters}
  \begin{tabular}{@{}llr@{}}
  \toprule
  Symbol & Meaning & Value\\
  \midrule
  $v_{\max}$ & speed limit; $V_s$ is the ball of this radius & $2$~m/s\\
  $a_{\max}$ & acceleration limit per axis (spans the window) & $2$~m/s$^2$\\
  $a_b$ & braking deceleration (admissibility), $a_b \le a_{\max}$ & $2$~m/s$^2$\\
  $\dt$ & control interval & $0.25$~s\\
  $\Delta v$ & resolution of the velocity grid & $0.25$~m/s\\
  $d_{\max}$ & cap on free distances (sensing range) & $3$~m\\
  $r$ & drone radius (obstacle inflation) & $0.2$~m\\
  $T$ & rollout horizon used for drawing predicted trajectories & $1$~s\\
  $(\alpha, \beta, \gamma)$ & weights of heading, clearance, velocity & $(0.6, 0.2, 0.2)$\\
  $\rho_{\mathrm{goal}}$ & goal tolerance & $0.2$~m\\
  \bottomrule
  \end{tabular}
\end{table}

% ---------------------------------------------------------------------
\section{The algorithm}
\label{sec:ch14-algorithm}

\Cref{alg:ch14-command} is one DWA step for the drone; \cref{alg:ch14-loop}
wraps it into the control loop. \Cref{tab:ch14-parameters} lists the
parameters.

\begin{algorithm}[htb]
\caption{One step of the dynamic window approach (drone form)}
\label{alg:ch14-command}
\KwIn{position $\pos$, current velocity $\vel_a$, target point $\vect{t}$, obstacles, parameters of \cref{tab:ch14-parameters}}
\KwOut{the velocity command for the next interval}
\Fn{\DwaCommand{$\pos$, $\vel_a$, $\vect{t}$, obstacles}}{
  $\vect{lo} \gets \max(\vel_a - a_{\max}\dt,\ -v_{\max})$; $\vect{hi} \gets \min(\vel_a + a_{\max}\dt,\ v_{\max})$\tcp*{window $V_d$, per axis}\label{alg:ch14-command:window}
  $C \gets$ grid points $\vel_a + \Delta v\,(i_1, \dots, i_{\dim})$, $\abs{i_k} \le a_{\max}\dt/\Delta v$, inside $[\vect{lo}, \vect{hi}]$ and with $\norm{\vel} \le v_{\max}$\label{alg:ch14-command:grid}\;
  $\vel_b \gets \vel_a \cdot \max\bigl(0,\ 1 - a_b\dt/\norm{\vel_a}\bigr)$; add $\vel_b$ to $C$\tcp*{braking candidate, always in $V_d$}\label{alg:ch14-command:brake}
  $G^* \gets -\infty$; $\vel^* \gets \vel_b$\;
  \ForEach{$\vel \in C$}{
    $d \gets$ \DwaFreeDistance{$\pos$, $\vel$, obstacles, $d_{\max}$}\tcp*{rollout: ray cast in direction $\vel$}\label{alg:ch14-command:rollout}
    \If{$\norm{\vel}\,\dt + \norm{\vel}^2/(2a_b) > d$}{\KwContinue\tcp*{not admissible, \cref{eq:ch14-vadm}}\label{alg:ch14-command:admissible}}
    $G \gets \alpha\,(1 - \theta(\vel)/\pi) + \beta\,\min(d, d_{\max})/d_{\max} + \gamma\,\norm{\vel}/v_{\max}$\label{alg:ch14-command:score}\;
    \If{$G > G^*$}{$G^* \gets G$; $\vel^* \gets \vel$\label{alg:ch14-command:argmax}\;}
  }
  \KwRet $\vel^*$\;
}
\end{algorithm}

\begin{algorithm}[htb]
\caption{The DWA control loop with a global path and hysteresis}
\label{alg:ch14-loop}
\KwIn{start $\pos_0$, goal $\pos_g$, obstacles, optional global path $P$ (a polyline), lookahead $L$, hysteresis margin $\eta \ge 0$}
\Fn{\DwaLoop{$\pos_0$, $\pos_g$, obstacles, $P$}}{
  $\pos \gets \pos_0$; $\vel_a \gets \vect{0}$\;
  \While{$\norm{\pos - \pos_g} > \rho_{\mathrm{goal}}$}{
    update the obstacle set from the sensors and the prediction layer\label{alg:ch14-loop:sense}\;
    $\vect{t} \gets \pos_g$ \textbf{if} $P = \emptyset$ \textbf{else} \DwaLookahead{$\pos$, $P$, $L$}\tcp*{heading reference}\label{alg:ch14-loop:target}
    $\vel \gets$ \DwaCommand{$\pos$, $\vel_a$, $\vect{t}$, obstacles}\;
    \If{$\eta > 0$ \KwAnd $\vel_a$ is admissible \KwAnd $G(\vel_a) \ge G(\vel) - \eta$}{$\vel \gets \vel_a$\tcp*{hysteresis: keep the old command}\label{alg:ch14-loop:hysteresis}}
    execute $\vel$ for $\dt$; $\pos \gets \pos + \vel\,\dt$; $\vel_a \gets \vel$\label{alg:ch14-loop:execute}\;
  }
}
\end{algorithm}

\paragraph{Walkthrough of \cref{alg:ch14-command}.}
Line~\ref{alg:ch14-command:window} builds the window of
\cref{def:ch14-window}, clipped to the speed limits so that the grid does
not leave $V_s$ along an axis; line~\ref{alg:ch14-command:grid} lays the
grid over it and drops the points outside the ball $V_s$. The grid is
centred on $\vel_a$, so ``keep the current velocity'' is always a candidate.
Line~\ref{alg:ch14-command:brake} adds the \textbf{braking
candidate}\index{braking candidate}\index{DWA!braking candidate}: same
direction as $\vel_a$, speed reduced by $a_b\dt$ (or zero). It lies in the
window because $a_b \le a_{\max}$, and \cref{thm:ch14-safety} shows that it
is always admissible when the previous command was; it is the reason why the
loop never runs out of safe options and it is also the command returned when
nothing else is admissible, which with static obstacles cannot happen.
Line~\ref{alg:ch14-command:rollout} is the rollout: for the drone a single
ray cast, for the differential drive a march along the arc of
\cref{eq:ch14-arc}. Line~\ref{alg:ch14-command:admissible} is the discrete
admissibility test \cref{eq:ch14-braking-discrete}; the code also applies it
to the distance to the goal, so that the drone brakes for the goal as it
would for a wall and does not orbit it. Line~\ref{alg:ch14-command:score}
evaluates \cref{eq:ch14-G} without smoothing, and
line~\ref{alg:ch14-command:argmax} keeps the best. Ties are broken by grid
order, which is deterministic but arbitrary; \cref{sec:ch14-properties}
returns to this.

\paragraph{Walkthrough of \cref{alg:ch14-loop}.}
Line~\ref{alg:ch14-loop:sense} is where DWA meets the rest of the system:
the obstacle set can be a map, a range scan, or the predicted positions of
other drones from \cref{ch:ch18,ch:ch20}. Line~\ref{alg:ch14-loop:target}
selects the target of the heading term. With $P = \emptyset$ it is the goal
itself, the pure ``greedy'' DWA of the original paper. With a global path it
is a \textbf{lookahead point}\index{lookahead point}: the point at arc length
$L$ beyond the point of $P$ closest to $\pos$, a carrot that the drone
chases along the path. Line~\ref{alg:ch14-loop:hysteresis} is the optional
hysteresis of \cref{sec:ch14-variants}. Line~\ref{alg:ch14-loop:execute}
moves the drone; in the simulator the command is executed exactly, on a real
drone a velocity controller tracks it and the window should be sized for
what that controller can actually deliver.

% ---------------------------------------------------------------------
\section{A worked example}
\label{sec:ch14-example}

\begin{example}[A drone approaching an obstacle]
\label{ex:ch14-approach}
A drone of radius $r = 0.2$~m is at $\pos = (0, 0)$, flying at
$\vel_a = (1.0, 0)$~m/s towards the goal $\pos_g = (5, 0)$. A disc obstacle
of radius $0.4$~m is centred at $(1.2, -0.1)$, almost on the line to the
goal and slightly below it. The parameters are those of
\cref{tab:ch14-parameters}. The instance is \code{worked\_example()} in
\code{ch14\_dwa.py}; every number in this section is printed by its self-test
and by \code{gen\_ch14\_traj.py}.
\end{example}

\paragraph{The window and the grid.}
With $a_{\max}\dt = 0.5$~m/s the window is
$V_d = [0.5, 1.5] \times [-0.5, 0.5]$, and with $\Delta v = 0.25$~m/s the
grid has $5 \times 5 = 25$ points, all inside the ball $\norm{\vel} \le 2$.
The braking candidate $(0.5, 0)$ is a grid point, so $C$ has $25$
elements (\cref{fig:ch14-example}b).

\paragraph{The admissibility test.}
The inflated obstacle has radius $0.6$~m. Straight ahead, the ray from the
origin along $x$ passes at distance $0.1$ from the centre and meets the disc
at $1.2 - \sqrt{0.6^2 - 0.1^2} = 0.608$~m. By \cref{eq:ch14-vadm} the
largest admissible speed in that direction is
$-0.5 + \sqrt{0.25 + 4 \cdot 0.608} = 1.138$~m/s: keeping $(1.0, 0)$ is
admissible (braking distance $0.25 + 0.25 = 0.5 \le 0.608$), but
$(1.25, 0)$ is not ($0.3125 + 0.39 = 0.70 > 0.608$). All ten candidates with
$v_x \ge 1.25$ except $(1.25, 0.5)$ are inadmissible: their rays hit the
disc within $0.6$ to $0.8$~m. Sixteen candidates survive; six of them, those
that steer clear of the disc at moderate speed, see no obstacle at all
within $d_{\max}$ and have $\dist = 3$.

\begin{figure}[tb]
  \centering
  \begin{subfigure}[b]{0.5\textwidth}\centering
    \inputfigure{ch14/example-scene}
    \caption{workspace}
  \end{subfigure}\hfill
  \begin{subfigure}[b]{0.48\textwidth}\centering
    \inputfigure{ch14/example-window}
    \caption{velocity space}
  \end{subfigure}
  \caption{\Cref{ex:ch14-approach}. (a) The drone (green circle) with the
  25 rollouts of one second: dotted red ones are inadmissible, blue ones
  admissible, the orange arrow is the chosen command $(1.0, 0.5)$; the dashed
  circle is the obstacle inflated by the drone radius, and the blue dots are
  the 13 positions of the closed-loop run to the goal. (b) The same step in
  velocity space: the dashed circle is $V_s$, the blue box the window $V_d$,
  the green region the admissible set $V_a$ of \cref{eq:ch14-vadm} (its bite
  points at the obstacle), the crosses the inadmissible grid points, the
  square the current velocity $\vel_a$ and the star the chosen command.}
  \label{fig:ch14-example}
\end{figure}

\paragraph{The scores.}
\Cref{tab:ch14-candidates} lists the three terms for 13 of the 25
candidates. Keeping $(1.0, 0)$ heads perfectly at the goal
($\mathrm{heading} = 1$) but has only $0.608/3 = 0.203$ of clearance, for
$G = 0.6 + 0.041 + 0.1 = 0.741$. The winner is $(1.0, 0.5)$: the ray misses
the inflated disc, so clearance is $1$, and although the heading drops to
$0.844$ (the rollout points $28^\circ$ away from the goal seen from the
predicted position $(0.25, 0.125)$) the total is $0.818$. The swerve is
towards $+y$, the side on which the obstacle leaves more room; the mirror
image $(1.0, -0.5)$ still hits the disc at $0.718$~m and scores only
$0.666$. Slower candidates with full clearance, $(0.75, 0.5)$ and
$(0.5, 0.25)$, lose on the velocity term. Note what decided the step: the
braking test removed every fast candidate that pointed at the obstacle, and
among the survivors the clearance term, not the heading term, chose the
side.

\begin{table}[tb]
  \centering\small
  \caption{Thirteen of the 25 candidates of \cref{ex:ch14-approach} with
  their free distance, admissibility and the three normalised terms;
  $G = 0.6\,\mathrm{heading} + 0.2\,\mathrm{clearance} + 0.2\,\mathrm{velocity}$.
  The chosen command is in bold.}
  \label{tab:ch14-candidates}
  \begin{tabular}{@{}rrrcrrrr@{}}
  \toprule
  $v_x$ & $v_y$ & $\dist$ [m] & admissible & heading & clearance & velocity & $G$\\
  \midrule
  1.50 & 0.00 & 0.608 & no & 1.000 & 0.203 & 0.750 & --\\
  1.25 & 0.00 & 0.608 & no & 1.000 & 0.203 & 0.625 & --\\
  1.00 & 0.00 & 0.608 & yes & 1.000 & 0.203 & 0.500 & 0.741\\
  0.50 & 0.00 & 0.608 & yes & 1.000 & 0.203 & 0.250 & 0.691\\
  1.50 & 0.50 & 0.739 & no & 0.889 & 0.246 & 0.791 & --\\
  1.25 & 0.50 & 0.812 & yes & 0.870 & 0.271 & 0.673 & 0.711\\
  \textbf{1.00} & \textbf{0.50} & 3.000 & yes & 0.844 & 1.000 & 0.559 & \textbf{0.818}\\
  0.75 & 0.50 & 3.000 & yes & 0.805 & 1.000 & 0.451 & 0.773\\
  0.50 & 0.50 & 3.000 & yes & 0.742 & 1.000 & 0.354 & 0.716\\
  0.50 & 0.25 & 3.000 & yes & 0.848 & 1.000 & 0.280 & 0.765\\
  1.00 & 0.25 & 0.682 & yes & 0.918 & 0.227 & 0.515 & 0.699\\
  1.00 & $-$0.50 & 0.718 & yes & 0.844 & 0.239 & 0.559 & 0.666\\
  0.50 & $-$0.50 & 3.000 & yes & 0.742 & 1.000 & 0.354 & 0.716\\
  \bottomrule
  \end{tabular}
\end{table}

\paragraph{The closed loop.}
\Cref{tab:ch14-trace} continues the run. In step~2 the window has moved to
$[0.5, 1.5] \times [0, 1]$ and the drone accelerates to $(1.5, 0.75)$, in
step~3 to $(1.75, 0.75)$: with the obstacle no longer in the way, the
velocity term takes over. From step~4 the heading term bends the flight back
towards the goal line, $(1.75, 0.25)$ and then $(1.75, -0.25)$ for six
steps, during which the heading rises to $0.99$. Steps~11--13 are the
approach to the goal: the goal-braking rule of
\cref{sec:ch14-algorithm} makes the fast candidates inadmissible and the
drone slows to $(1.25, -0.25)$, $(1.0, -0.25)$ and finally the braking
candidate $(0.51, -0.13)$, which is not a grid point. After 13 steps
($3.25$~s) the drone is at $(4.82, 0.03)$, inside the goal tolerance. The
path is $4.98$~m long, and the smallest gap between the drone's hull and
the obstacle along the way was $1.6$~cm: DWA with a dominant heading term
shaves obstacles, because the clearance term only counts what lies ahead on
the rollout, not what lies beside it. Raising $\beta$, inflating the obstacles
by a safety margin, or the smoothing $\sigma$ all widen the berth.

\begin{table}[tb]
  \centering\small
  \caption{Trace of the closed-loop run of \cref{ex:ch14-approach}: position
  and velocity at the start of each step, the chosen command and its terms.
  Every chosen command has $\dist = 3$~m (no obstacle on its ray within
  $d_{\max}$), so clearance is $1$ throughout.}
  \label{tab:ch14-trace}
  \begin{tabular}{@{}rlllrrr@{}}
  \toprule
  Step & $\pos$ & $\vel_a$ & chosen $\vel$ & heading & velocity & $G$\\
  \midrule
  1 & $(0.00, 0.00)$ & $(1.00, 0.00)$ & $(1.00, 0.50)$ & 0.844 & 0.559 & 0.818\\
  2 & $(0.25, 0.12)$ & $(1.00, 0.50)$ & $(1.50, 0.75)$ & 0.830 & 0.839 & 0.866\\
  3 & $(0.62, 0.31)$ & $(1.50, 0.75)$ & $(1.75, 0.75)$ & 0.831 & 0.952 & 0.889\\
  4 & $(1.06, 0.50)$ & $(1.75, 0.75)$ & $(1.75, 0.25)$ & 0.904 & 0.884 & 0.919\\
  5 & $(1.50, 0.56)$ & $(1.75, 0.25)$ & $(1.75, -0.25)$ & 0.994 & 0.884 & 0.973\\
  6 & $(1.94, 0.50)$ & $(1.75, -0.25)$ & $(1.75, -0.25)$ & 0.993 & 0.884 & 0.972\\
  7 & $(2.38, 0.44)$ & $(1.75, -0.25)$ & $(1.75, -0.25)$ & 0.991 & 0.884 & 0.971\\
  8 & $(2.81, 0.38)$ & $(1.75, -0.25)$ & $(1.75, -0.25)$ & 0.989 & 0.884 & 0.970\\
  9 & $(3.25, 0.31)$ & $(1.75, -0.25)$ & $(1.75, -0.25)$ & 0.985 & 0.884 & 0.968\\
  10 & $(3.69, 0.25)$ & $(1.75, -0.25)$ & $(1.75, -0.25)$ & 0.978 & 0.884 & 0.964\\
  11 & $(4.12, 0.19)$ & $(1.75, -0.25)$ & $(1.25, -0.25)$ & 0.993 & 0.637 & 0.923\\
  12 & $(4.44, 0.12)$ & $(1.25, -0.25)$ & $(1.00, -0.25)$ & 0.985 & 0.515 & 0.894\\
  13 & $(4.69, 0.06)$ & $(1.00, -0.25)$ & $(0.51, -0.13)$ & 0.974 & 0.265 & 0.838\\
  \bottomrule
  \end{tabular}
